\section{The Data Kernel Perspective Space}
Before we present our results related to statistical inference on black-box generative models, we first describe how to obtain finite-dimensional vector representations of the models. 
% Vector representations of objects facilitates statistical analysis of the objects with conventional methods. 
The particular method for obtaining vector representations for generative models that we study herein has previously been discussed in \cite{helm2024tracking} and \cite{aranyak}.

% Representing objects of interests with finite-dimensional vectors often facilitates statistical analysis of the objects with conventional statistical tools. 
% For instance, one can group a given set of objects into different clusters by clustering their vector representations. 
% Before we present our key results concerning statistical inference on vector representations of generative models, we describe here how we obtain the vector representations of the generative models. 
% We remind the reader here that this particular method is not novel to this paper, it has already been discussed in two previous works \cite{duderstadt2023comparing} and \cite{aranyak}.

% -- such as predicting if a model has had access to sensitive information or predicting a model's performance on a benchmark. 

% \subsection*{Methodology}
Let $ f: \mathcal{Q} \to \mathcal{X} $ be a generative model from a query space $ \mathcal{Q} $ to an output space $ \mathcal{X} $. 
Given $ q \in \mathcal{Q} $, we assume that $ f(q) $ produces a random element in $ \mathcal{X} $.
In our setting there are $ n $ models $f_{1},\hdots f_{n} $ and $ m $ queries $ q_{1}, \hdots, q_{m} $ with $ q_{j} \in \mathcal{Q} $.
We assume that each model responds to every query $ r $ independent times and let $ f_{i}(q_{j})_{k} $ be the $ k $-th response to $ q_{j} $ from $ f_{i} $. 
% Throughout our discussion we let $ i\;\text{\&} \;i'$ index models, $ j \;\text{\&} \;j' $ index queries, and $ k $ index response replicates.



% Suppose there are $n$ generative models $f_1,\dots f_n$ and $m$ queries $q_1,\dots q_m$.  and let $f_i(q_j)_k$ denote the $k$-th replicate of the response of $f_i$ to $q_j$. For every $i,j$, the replicates $f_i(q_j)_1,\dots f_i(q_j)_r$ are independently and identically distributed over some pool of responses. 


% The setting is as follows. Suppose there are $n$ generative models $f_1,\dots f_n$ and $m$ queries $q_1,\dots q_m$. Each model $f_i$ responds to every query $q_j$ on $r$ replicates, and let $f_i(q_j)_k$ denote the $k$-th replicate of the response of $f_i$ to $q_j$. For every $i,j$, the replicates $f_i(q_j)_1,\dots f_i(q_j)_r$ are independently and identically distributed over some pool of responses. 


We let $ g: \mathcal{X} \to \mathbb{R}^{p} $ be an embedding function that maps a model response to a $p$-dimensional real-valued vector.
Further, we let $x_{ijk}=g(f_i(q_j)_k)$ denote the embedding of $ f_{i}(q_{j})_{k} $ and $ F_{ij} $ denote the distribution of $ x_{ijk} $ on $ \mathbb{R}^{p} $. 
That is, $x_{ij1},\dots x_{ijr} \iid  F_{ij} $. 
% \footnote{We sometimes refer to $ F_{ij} $ as the \textit{generation distribution} of $ f_{i} $ with respect to $ q_{j} $.}

We let $ \bar{x}_{ij} = \frac{1}{r} \sum_{k=1}^{r} x_{ijk} $ be the empirical average embedded response of $f_{i}$ to $q_{j}$ and $ \bar{X}_{i} \in \mathbb{R}^{m \times p}$ denote the matrix whose $ j $-th row is $ \bar{x}_{ij} $.
We view $ \bar{X}_{i} $ as a matrix representation of model $ f_{i} $ with respect to $ \{ q_1,\dots, q_m \}$ and $ g $. 
We define $ D $ as the $ n \times n $ pairwise distance matrix with entries 
\begin{align}
\label{euc-dist-matrix}
    D_{ii'} = \frac{1}{m} \big|\big| \bar{X}_{i} - \bar{X}_{i'} \big|\big|_{F}
\end{align} and note that $ D_{ii'} $ captures the difference in empirical average model behavior between $ f_{i} $ and $ f_{i'} $ with respect to $ \{q_{1}, \hdots, q_{m}\} $ and $ g $.

With the form described by Eq. \eqref{euc-dist-matrix}, $ D $ is a rescaled Euclidean distance matrix. 
Hence, multidimensional scaling (MDS) of $ D $ yields $ d $-dimensional Euclidean representations of the matrices $ \bar{X}_{i} $ (and thereby of the models $f_i$) with respect to $ \{q_{1}, \hdots, q_{m}\} $ and $ g $ \citep{torgerson1952multidimensional}.
Letting $ \widehat{\psi} := \text{MDS}(D) \in \mathbb{R}^{n \times d} $, we refer to $ \widehat{\psi} $ as the \textit{data kernel perspective space} (DKPS).
We emphasize that the $ i $-th row of the DKPS -- $ \widehat{\psi}_{i} $ -- is a $ d $-dimensional vector representation of the generative model $ f_{i} $.

% In particular, let $ D^{2} $ be the $ n \times n $ matrix with entries $ D_{ii'}^{2} $, $ e_{n} \in \mathbb{R}^{n}$ be the vector of all ones, $ I_{n} $ be the $ n \times n $ identity matrix, and $ P = I_{n} - e_{n}e_{n}^{T} / n $. Then, letting $ B = -\frac{1}{2} P D^{2} P^{T} $ and $ v_{\ell} $ be the eigenvector associated with the $ \ell $th largest eigenvalue $ \lambda_{\ell} $ of $ B $, the rows of the matrix $ \widehat{\psi} \in \mathbb{R}^{n \times d} $ with columns $ \lambda_{\ell}^{2} v_{\ell} $ are $ d $-dimensional Euclidean representations of the models with respect to $ Q $. 



% \subsubsection*{Properties of the Euclidean representations}

% \clearpage
\subsection{Analytical properties of the DKPS}

While the DKPS representations of the models are Euclidean objects, it is not possible to comment on their properties without imposing constraints on the $ F_{ij} $.
%To facilitate analysis we assume each query is from the query distribution $ G $ with support $ \mathscr{Q} $  
% let the distribution of response distributions $ F_{ij} $ with $ q_{j} \sim G $ be $ \mathcal{F}_{i}(G) $ for fixed $ f_{i} $ (i.e., $ F_{ij} \sim \mathcal{F}_{i}(G) $) 
We let $ \mu_i \in \mathbb{R}^{m \times p}$ be the population counterpart of $ \bar{X}_i $ whose $ j $-th row is $ \mathbb{E}_{x \sim F_{ij}}(x)$. 
Similarly, we let $ \Delta $ be the population counterpart of $ D $ whose $(i,i')$-th element is $ \Delta_{ii'} = \frac{1}{m}||\mu_{i} - \mu_{i'}||_{F} $.
We assume
 $ \Delta^{*}_{i i'} = \lim \frac{1}{m} || \mu_{i} - \mu_{i'}||_{F} $ for every pair $ (i,i') $ as $ m, r \to \infty $ exists
% In our setting, where both $ m, r \to \infty $, we assume
 % $ \Delta^{*}_{i i'} = \lim_{m,r \to \infty} \frac{1}{m} || \mu_{i} - \mu_{i'}||_{F} $ for every pair $ (i,i') $.
% \begin{align*}
% \mu_{i}(G)~:=~\mathbb{E}_{F_{ij}\sim\mathcal{F}(G)}\left[\mathbb{E}_{x\sim F_{ij}}\left[g(x))\right]\right].
% \end{align*}
and that $ g $ is bounded.
% \footnote{Boundedness is typically satisfied in practice for well-defined $ g $. For language models, an element of $ \mathcal{X} $ is a finite sequence from a finite vocabulary; for text-to-image models $ \mathcal{X} = \{0, \hdots, 255\}^{3}$ is finite.}

% . For example, for language models elements of both $ \mathscr{Q} $ and $ \mathcal{X} $ are sequences of tokens from a finite vocabulary where the length of an element in $ \mathscr{Q} $ is bounded above by the size of the context window and the length of an element of $ \mc{X} $ is bounded above by a user-defined maximum sequence length.

%( i.e., $ |\mathscr{Q}| < \infty $ and $ |\mathcal{X}| < \infty $).  
% Together, these imply that the image of $ \mathcal{X} $ under $ g $ is bounded and, hence, that $ g(x) $ with $ x \sim F_{ij} $ has finite moments and $ \mathbb{E}_{F_{ij} \sim \mathcal{F}(G)}\left[\mathbb{E}_{x \sim F_{ij}}\left[g(x)\right] \right] $ is finite. 

In our setting, under appropriate assumptions described in \citep{aranyak},
% and letting 
% \begin{align*}
% \mu_{i}(G)~:=~\mathbb{E}_{F_{ij}\sim\mathcal{F}(G)}\left[\mathbb{E}_{x\sim F_{ij}}\left[g(x))\right]\right],
% \end{align*}
%we have
% $|D_{i i'}-\Delta_{i i'}(G)| \to 0$  as $m,r \to \infty$ 
% with high probability for all $ i, i'$.
% where $\Delta_{i i'}(G)=\frac{1}{m}
% \left\lVert 
% \mu_i(G)-\mu_{i'}(G)
% \right\rVert
% $. 
% Additionally, we assume 
% $\lim_{m,r \to \infty}
% \Delta_{i i'}(G)
% \to \Delta^{(\infty)}(G)
% $  for all $n$, where $\Delta^{(\infty)}
% \in \mathbb{R}^{n \times n}
% $ is a dissimilarity matrix, so we have
% for all $n$,
\begin{equation}
    D_{i i'}=
    \frac{1}{m}
    \left\lVert 
\bar{X}_i-\bar{X}_{i'}
    \right\rVert
    \to 
    \Delta_{i i'}^{*}
\end{equation}
with high probability
as $m,r \to \infty$ for all $ (i, i') \in \{1, \hdots, n\} \times \{1, \hdots, n\} $.
%\begin{align}
%\label{eq:entry-of-D2}
%    D_{ii'} = \frac{1}{m} \big\| \big(\bar{X}_{i} - \bar{X}_{i'}\big) \big\|_{F} 
%    \xrightarrow[]{}
%    \big\|  \mu_{i}(G) - \mu_{i'}(G) \big\|_{F} =: \Delta_{ii'}(G)
%\end{align} 
%as $ R, m \to \infty $ for all $ (i, i') \in \{1, \hdots, n\} \times \{1, \hdots, n\} $ \citep{aranyak}.  
Further, there exists $ \psi := \text{MDS}(\Delta^{*}) \in \mathbb{R}^{m \times d} $ such that $ \widehat{\psi} \to \psi $. 
That is, the configuration $ \psi $ is the population counterpart of $ \widehat{\psi} $. $ \psi $ captures the true geometry of the mean discrepancies of the model responses with respect to the queries  $ \{q_1,\hdots, q_m \} $.
%We drop the argument of $ \psi $ when the distribution of the queries is clear.
% and hinges on the boundedness of $ g(x) $. 
Importantly, in settings where the models are not fixed and are instead assumed to be $ i.i.d. $ realizations from a model distribution $ P_{model} $, the distance matrix $ D $ converges to $ \Delta^{*} $ and, under technical assumptions, there exists $ \psi $ such that $ \widehat{\psi} \to \psi $ as $ m, r \to \infty $ for all $n$ \citep{aranyak}.

% In settings where the query space contains only a single element or where each $ F_{ij} $ is a point mass on $ x_{ij} $, $ \Delta_{ii'} $ is exactly the maximum mean discrepancy \citep{gretton2012kernel} of the distributions $ F_{ij} $ and $ F_{i'j} $ under a linear map applied to $ g(x) $. 
% For settings where $ | \mathscr{Q} | > 1 $ and the $ F_{ij} $ are not all point masses, the analysis is more complicated since both $ F_{ij} $ and $ \bar{x}_{ij}$ 
% are random variables. 
% The entry-wise convergence of $ D $ to $ \Delta(G) $ implies the convergence of $ Z_{i} $ to an analogously defined $ Z^{*}_{i}(G) \in \mathbb{R}^{p} $ up to an orthogonal transformation. The proof of this is provided in Appendix [ref].


% While we do not explore alternative dissimilarities on the $ m \times p \times R $ matrices obtained after embedding the $ R $ samples from each $ F_{ij} $ in this paper, we note that it is possible (and perhaps worthwhile) to consider more expressive notion of dissimilarity on the models. 
% For example, we could replace the row-wise difference with the average of a row-wise MMD with a universal kernel to ensure that we capture all relevant differences between the response distributions induced by different $ f_{i} $.
% The resulting $ n \times n $ pairwise distance matrix will likely not be a Euclidean distance matrix, though transforming it into useful and analyzable representations of the models is sometimes possible [cite, Trosset?].

% \clearpage

\subsection{Statistical inference in the DKPS}
\label{subsec:inference}

We now extend the consistency of $ \widehat{\psi} $ to $ \psi $ to model-level inference.
Consider the classical statistical learning problem \citep[Chapter~2]{hastie2009elements} in the context of a collection of generative models: given training data $ \mathcal{T}_{n} = \{ (f_{1}, y_{1}), \hdots, (f_{n}, y_{n}) \} $ assumed to be  $ i.i.d. $ realizations from the joint distribution $ P_{f Y} $, choose the decision function $ h: \mathcal{F} \to \mathcal{Y} $ that minimizes the expected value of a loss function $ \ell $ with respect to $ P_{f Y} $ within a class of decision functions $ \mathcal{H} $ for a test observation $ f $ assumed to be an independent realization from the marginal distribution $ P_{f}$. 
Or, with $ \mathcal{R}_{\ell}(P_{f Y}, h) := \mathbb{E}_{P_{f Y}}\left[\ell(h(f), y)\right] $, select $ h^* $ such that 
\begin{align*}
\label{eq:risk}
    h^* \in \argmin_{h \in \mathcal{H}} \mathcal{R}_{\ell}(P_{f Y}, h). 
\end{align*}
We let $ \mathcal{R}_{\ell}^{*}(P_{f Y}, \mathcal{H})=
\mathcal{R}_{\ell}(P_{f Y}, h^*)
$ denote the expected loss (or ``risk") of $ h^* $. 

The joint distribution $ P_{f Y} $ is often unavailable and the decision function is selected based on the training data. 
We let $ h(\; \cdot \;; \mathcal{T}_{n}) $ denote such a decision function and say \textit{the sequence of decision functions $ \left(h(\; \cdot \;; \mathcal{T}_{1}), \hdots, h(\;\cdot\;; \mathcal{T}_{n})\right) $ is consistent for $ P_{f Y} $ with respect to $ \mathcal{H}$} if \begin{align*}
Pr\big(\;\big|  \mathcal{R}_{\ell}(P_{f Y}, h(\;\cdot\;; \mathcal{T}_{n})) - \mathcal{R}^{*}_{\ell}(P_{f Y}, \mathcal{H}) \big| > \epsilon\big) \to 0
\end{align*} as $ n \to \infty $.

% Consider the classical statistical learning problem \citep[Chapter~2]{hastie2009elements} in the context of a collection of generative models: given training data $ \mathcal{T}_{n} = \{ (f_{1}, y_{1}), \hdots, (f_{n}, y_{n}) \} $ assumed to be  $ i.i.d. $ realizations from the joint distribution $ P_{\psi Y} $, choose the decision function $ h: \mathcal{X} \to \mathcal{Y} $ that minimizes the expected value of a loss function $ \ell $ with respect to $ P_{\psi Y} $ within a class of decision functions $ \mathcal{H} $ for a test observation $ \psi $ assumed to be an independent realization from the marginal distribution $ P_{\psi}$. Or, with $ \mathcal{R}_{\ell}(P_{\psi Y}, h) := \mathbb{E}_{P_{\psi Y}}\left[\ell(h(\psi), y)\right] $, select $ h^* $ such that 
% \begin{align*}
% \label{eq:risk}
%     h^* \in \argmin_{h \in \mathcal{H}} \mathcal{R}_{\ell}(P_{\psi Y}, h). 
% \end{align*}
% We let $ \mathcal{R}_{\ell}^{*}(P_{\psi Y}, \mathcal{H})=
% \mathcal{R}_{\ell}(P_{\psi Y}, h^*)
% $ denote the expected loss of $ h^* $.

% The joint distribution $ P_{\psi Y} $ is often unavailable and the decision function is selected based on $ \mathcal{T}_{n} $. 
% We let $ h(\; \cdot \;; \mathcal{T}_{n}) $ denote such a decision function and say \textit{the sequence of decision functions $ \left(h(\; \cdot \;; \mathcal{T}_{1}), \hdots, h(\;\cdot\;; \mathcal{T}_{n})\right) $ is consistent for $ P_{\psi Y} $ with respect to $ \mathcal{H}$} if \begin{align*}
% Pr\big(\;\big|  \mathcal{R}_{\ell}(P_{\psi Y}, h(\;\cdot\;; \mathcal{T}_{n})) - \mathcal{R}^{*}_{\ell}(P_{\psi Y}, \mathcal{H}) \big| > \epsilon\big) \to 0
% \end{align*} as $ n \to \infty $.

In the black-box setting we do not have direct access to useful representations of the models. 
Instead, we must query each model repeatedly and use the representations $ \{ \psi_1,\hdots, \psi_n \} $ as proxies for the models $ \{ f_{1}, \hdots, f_{n}\} $. 
Hence, $\mathcal{T}_n=\{ (\psi_1,y_1),\hdots ,(\psi_n,y_n) \}$
% With a slight abuse of notation, we henceforth denote the training data by $\mathcal{T}_n=\{ (\psi_1,y_1),\hdots ,(\psi_n,y_n) \}$ 
where $(\psi_i,y_i) \sim^{iid} P_{\psi Y}$ and our goal is to choose a decision function $h^*$ which minimizes the average loss 
$ \mathcal{R}_{\ell}(P_{\psi Y}, h) := \mathbb{E}_{P_{\psi Y}}\left[\ell(h(\psi), y)\right] $ within an appropriately defined $\mathcal{H}$.
Note that the true $ \{\psi_{1}, \hdots, \psi_{n}\} $ are not known \textit{a priori} and must be estimated via $ \{\widehat{\psi}_{1}, \hdots, \widehat{\psi}_{n}\} $.
We let $ \widehat{\mathcal{T}}_{n} $ be the training data where $ \psi_{i} $ is replaced with $ \widehat{\psi}_{i} $.
% Note that in order to compute $\psi_i$ we need access to the response distributions $F_{ij} $, hence we estimate $\psi_i$ using its sample counterpart $\widehat{\psi}_i$. In the abovementioned statistical learning problem, instead of dealing with the dataset
%Instead, we have access to representations of their generative properties with respect to $ \{q_{1}, \hdots, q_{m}\} $ captured via the data kernel perspective space.
%That is, we let $ \mathcal{T}_{n} = \{\psi_{1}, y_{1}), \hdots, \psi_{n}, y_{n})\} $ and the corresponding joint distribution, class of decision function, and loss function defined on $ \psi $ parameterize the inference problem.
% and the test observation $\psi_{n+1} $.
%Further, we do not observe $ \psi_{i} $ directly and must estimate it via $ \widehat{\psi}_{i} $.
% We note that the $ \widehat{\psi}_{i} $'s are errorful observations of the true-but-unknown $ \psi_{i} $. 
% Letting $ \{(f_{1}, y_{1}), \hdots, (f_{n}, y_{n})\} $ be $ i.i.d. $ realizations from a distribution on generative models and $ P_{\psi Y} $ be the corresponding distribution induced by MDS($\Delta$),

Two important extensions of the DKPS consistency results in the context of the classical statistical learning problem are the following theorems:
\begin{theorem}
\label{thm:rule-convergence}
    Under technical assumptions described in Appendix \ref{app:proofs}, \begin{align*}
    \mathcal{R}_{\ell}(P_{\psi Y}, h(\; \cdot \;; \widehat{\mathcal{T}}_{n})) \to \mathcal{R}_{\ell}(P_{\psi Y}, h(\; \cdot \;; \mathcal{T}_{n})) 
    \end{align*} as $ m, r \to \infty $, for every $n$.
\end{theorem} That is, the risk of a decision function based on $\widehat{\mathcal{T}}_n=
\{
(\widehat{\psi}_1,y_1),\hdots ,
(\widehat{\psi}_n,y_n)
\}
$ converges to the risk of the decision function based on the true-but-unknown $ \mathcal{T}_{n} $.
For situations where $ \{\psi_{1}, \hdots, \psi_{n}\} $ are good proxies, Theorem \ref{thm:rule-convergence} states that increasing the number of queries and/or number of replicates per query will improve the performance of decision functions trained on $ \widehat{\mathcal{T}}_{n} $.
% and
\begin{theorem}
\label{thm:consistency}
    Under technical assumptions described in Appendix \ref{app:proofs}, if $ (h(\; \cdot \;; \mathcal{T}_{1}), \hdots, h(\; \cdot \;; \mathcal{T}_{n})) $ is consistent for $ P_{\psi Y} $ with respect to
    $ \mathcal{H} $, then $ (h(\; \cdot \;; \widehat{\mathcal{T}}_{1}), \hdots, h(\; \cdot \;; \widehat{\mathcal{T}}_{n})) $ is consistent for 
    $ P_{\psi Y} $ 
    with respect to
    $ \mathcal{H} $ as $ n, m, r \to \infty $.
\end{theorem}
Theorem \ref{thm:consistency} states that if the sequence of decision functions learned from $ \mathcal{T}_{n} $ is consistent then the sequence of decision functions learned from the analogous $ \widehat{\mathcal{T}}_{n} $ is also consistent.
This result suggests that inference performance should improve with more training data.
While a logical follow-on to Theorem \ref{thm:rule-convergence} and a well-understood principle in classical statistical learning settings, the consistency of the DKPS in the setting where the number of models grows is a non-trivial extension of the fixed $ n $ case.

The proofs of Theorems \ref{thm:rule-convergence} and \ref{thm:consistency} are provided in Appendix \ref{app:proofs}.

\subsubsection{Non-standard considerations}
\label{subsec:non-standard-considerations}

Given that our analysis is based on estimates $ \{\widehat{\psi}_{1}, \hdots, \widehat{\psi}_{n} \}$ of proxies $ \{\psi_{1}, \hdots, \psi_{n}\} $ of the objects of interest $ \{f_{1}, \hdots, f_{n}\}$, it is important to note some non-standard theoretical considerations.
We let $ \mathcal{R}^{*}_{\ell}(P_{fY}) := \inf_{\{h: \mathcal{F} \to \mathcal{Y}\}}\mathcal{R}_{\ell}(P_{fY}, h) $ be the Bayes optimal risk of $ P_{fY} $ and $ P_{query} $ be a probability distribution on queries. 

For example, while our theoretical results provide insights as to how performance will be affected by increasing $ n, m, $ and $ r $, our results do not comment on the magnitude of $ | \mathcal{R}^{*}_{\ell}(P_{\psi Y}) - \mathcal{R}^{*}_{\ell}(P_{fY}) | $, which measures how good of a proxy $ \{\psi_{1}, \hdots, \psi_{n}\} $ is for $ \{f_{1}, \hdots, f_{n}\} $.
In particular, if $ | \mathcal{R}^{*}_{\ell}(P_{\psi Y}) - \mathcal{R}^{*}_{\ell}(P_{fY}) | = 0 $ then the $ \psi_{i} $ are lossless proxies of the $ f_{i} $ with respect to $ y $ \citep[Chapter~32]{devroye2013probabilistic}.
We expect it to be challenging to understand $ | \mathcal{R}^{*}_{\ell}(P_{\psi Y}) - \mathcal{R}^{*}_{\ell}(P_{fY}) | $ in a general setting.

Similarly, let $ q_{1}, \hdots, q_{m} \sim^{iid} P_{query} $ and $ q'_{1}, \hdots, q'_{m} \sim^{iid} P'_{query} $ with $ P_{query} \neq P'_{query} $.
The proxies of the models induced by these two query sets -- $ \{\psi_{1}, \hdots, \psi_{n}\} $ and $ \{\psi_{1}', \hdots, \psi_{n}'\} $, respectively -- will be of different quality.
Without an initial understanding of the models and their relationship with the covariates, we expect that it will be impossible to know which query distribution is preferred \textit{a priori}. 
We highlight the importance of $ P_{query} $ in an experimental setting below. 

Lastly, we note that representations of the models that we analyze herein capture the intrinsic geometry of the mean discrepancies of the models for the queries $ \{q_{1}, \hdots, q_{n}\} $.
For covariates that cannot be described as a function of the mean discrepancies of the models, it is unclear how to interpret Theorems \ref{thm:rule-convergence} and \ref{thm:consistency}.
We discuss potential alternatives to the representations that we study in the limitations section.

\subsection{An illustrative example -- ``Was RA Fisher great?"}
\label{subsec:fisher}
\begin{figure*}[t!]
    \centering
    \begin{subfigure}[t]{0.45\textwidth}
        \centering
        \includegraphics[width=\textwidth]{figures/was-ra-fisher-great/3d-surface-plot-091524.png}
    \end{subfigure}
    ~ 
    \begin{subfigure}[t]{0.45\textwidth}
        \centering
        \includegraphics[width=\textwidth]{figures/was-ra-fisher-great/ideal_fig1_nmr.png}
    \end{subfigure}
    \caption{
    \textbf{Left.} The 2-d Data Kernel Perspective Space (DKPS) and covariate surface for a collection of 550 models parameterized by fixed augmentations.
    \textbf{Right.} The performance of the 1-nearest neighbor regressor in DKPS for predicting the probability that an unlabeled model responds ``yes" to ``Was RA Fisher great?". 
    % The first dimension of DKPS correlates strongly with the covariate. Performance of the 1-NN regressor depends on both the number of models with known covariates and the number of queries used to induce the DKPS.
    }
    \label{fig:ra-fisher}
\end{figure*}

The first model-level inference task we study is a toy example where the task is to predict the probability that a model will respond ``yes" to the question ``Was RA Fisher great?". 
The question is chosen due to its subjectiveness -- there is no correct answer to the question of a person's greatness -- and its duality -- Ronald A. Fisher is considered one of the most influential statisticians in history and is considered an advocate of the 20th century Eugenics movement \citep{bodmer2021outstanding}. 

We use a 4-bit version of Meta's \texttt{LLaMA-2-7B-Chat} \citep{touvron2023llama} as a base model and consider a collection of models parameterized by fixed context augmentations. 
Each augmentation $ a_{i} $ contains information related to RA Fisher's statistical achievements (i.e., $ a_{i} $ = ``RA Fisher pioneered the principles of the design of experiments") or to his involvement in the eugenics movement (i.e., $ a_{i} $ = ``RA Fisher's view on eugenics were primarily based on anecdotes and prejudice.") and is prepended to every query.
The covariate corresponding to a given model is calculated by prompting the base model with the appropriately formatted prompt ``Give a precise answer to the question based on the context. Don't be verbose. The answer should be either a yes or a no. $ a_{i} $. Was RA Fisher great?" until there are 100 valid responses. We let $ y_{i} $ be the average number of ``yes"es.

To induce a DKPS for this task, we consider queries sampled from OpenAI's ChatGPT with the prompt ``Provide 100 questions related to RA Fisher.".  
For a given query $ q_{j} $ we prompt the base model with the appropriately formatted prompt ``$a_{i} \; q_{j}$" and fix $ r = 1 $.
The left figure of Figure \ref{fig:ra-fisher} is a 3-d figure where the first two dimensions are the DKPS of the $ n = 550 $ ($ 275 $ statistics augmentations, $ 275 $ eugenics augmentations) models induced with $ m = 100 $ queries. Model responses are embedded by averaging the per-token last layer activation of the base model.  
The third dimension is an interpolated $ y_{i} $ surface with a linear kernel \citep{du2008radial}. 
We emphasize that the dots populating the left figure of Figure \ref{fig:ra-fisher} are representations of the models with respect to all 100 queries -- not just vector representations of a query or a response to a query -- and that the $ y $-surface corresponds to the average response of each model to a query not included in the 100 queries.

The first dimension of the DKPS is clearly capable of distinguishing between models adorned with ``statistics" augmentations and models adorned with ``eugenics" augmentations. 
The shape of the interpolated covariate surface is highly correlated with this feature.
A description of the augmentations that parameterize the models and the queries used to induce the DKPS is provided in Appendix \ref{app:fisher}.

The right figure of Figure \ref{fig:ra-fisher} shows the performance of a 1-nearest neighbor (1-NN) regressor in DKPS for a varying number of labeled models and a varying number of queries. 
The DKPS is induced with $ n $ models and the regressor is trained with $ n - 1 $ of the model-level covariates. 
The mean squared error reported is the error of the 1-NN regressor for predicting the ``left out" model's covariate ($\pm$ 1 S.E.).
As Theorems \ref{thm:rule-convergence} and \ref{thm:consistency} suggest, the performance of the regressor is dependent on both the number of models and the number of queries: the more models and the more queries the better.
The scale of the impact of more models and more queries, however, depends on its counterpart. 
The amount of models does not have a large impact on predictive performance if the number of queries is small. 
The amount of queries has a large impact on performance regardless of the number of models.